\documentclass[11pt,reqno]{amsart}

\usepackage[margin=1in]{geometry}
\usepackage[T1]{fontenc}
\usepackage{lmodern}
\usepackage{microtype}
\usepackage{mathtools}
\usepackage{amssymb}
\usepackage{braket}
\usepackage{xcolor}
\usepackage[colorlinks=true,linkcolor=blue!55!black,citecolor=blue!55!black,
  urlcolor=blue!55!black]{hyperref}

\theoremstyle{plain}
\newtheorem{theorem}{Theorem}[section]
\newtheorem{corollary}[theorem]{Corollary}
\newtheorem{proposition}[theorem]{Proposition}
\newtheorem{lemma}[theorem]{Lemma}
\theoremstyle{remark}
\newtheorem{remark}[theorem]{Remark}

\usepackage[nameinlink,capitalize,noabbrev]{cleveref}

\DeclareMathOperator{\Tr}{Tr}
\DeclareMathOperator{\rank}{rank}
\DeclareMathOperator{\id}{id}
\DeclareMathOperator{\vecop}{vec}
\newcommand{\cL}{\mathcal L}
\newcommand{\cS}{\mathcal S}
\newcommand{\cU}{\mathcal U}
\newcommand{\cD}{\mathcal D}
\newcommand{\ketbra}[2]{\lvert #1\rangle\!\langle #2\rvert}
\newcommand{\ip}[2]{\langle #1,#2\rangle}
\newcommand{\norm}[1]{\lVert #1\rVert}
\newcommand{\abs}[1]{\lvert #1\rvert}
\newcommand{\C}{\mathbb C}

\title{Exact finite-level free-spectrahedral inclusion\\
from a two-layer graph operator system}
\author{Elazar Gershuni}
\date{\today}

\begin{document}
\begin{abstract}
Pair amplification turns a Schmidt-rank-two estimate on a Kraus dictionary into
$R$-positivity of a corrected map, hence into an inclusion of free spectrahedra
valid through matrix level $R$ and failing at level $R+1$.  Read on a full
matrix dictionary the resulting statement is tomographically complete, so it
certifies nothing a Choi computation would not.  We exhibit a family of
edge-indexed Kraus maps $\Phi_G$ on $M_{2m}$, one for each simple graph $G$ on
$m$ vertices, with $\beta_2(\Phi_G)=1$ exactly, and we show that the amplified
coefficient $R-1$ is attained whenever $G$ contains an $R$-clique.  The
coefficient system can then be compressed: not to the graph operator system
$\cS_G$, whose one-corner compression destroys the protected direction, but to
the two-layer system $\cU_G=\cS_G\otimes\cD_2$, which is invariant, carries
$2m(1+k)-1$ traceless coefficients for a $k$-regular $G$, and still contains a
positive matrix-level element witnessing sharpness.  For the rook graph
$H(2,q)$ this gives exact inclusion through every level $R<q$ using
$2q^2(2q-1)-1$ coefficients in place of $4q^4-1$.
\end{abstract}
\maketitle

\section{Setting}

Let $\Phi(X)=\sum_aA_aXA_a^*$ be completely positive on $M_d$ with every Kraus
operator transpose-symmetric, $A_a^T=A_a$.  Write
\[
    2\beta_2(\Phi)
    =\sup_{\norm c_2=1}\bigl(s_1(K_c)^2+s_2(K_c)^2\bigr),
    \qquad K_c=\sum_ac_aA_a,
\]
and, for an integer $R\ge1$,
\begin{equation}
\label{eq:theta}
    \Theta^\Phi_{R,\lambda}
    :=\Phi\circ\tau+\lambda\Bigl(\Delta_d-\tfrac1R\id\Bigr),
    \qquad
    \Delta_d(X)=\Tr(X)I_d,
\end{equation}
with $\tau$ the transpose.  Pair amplification (\texttt{rank-r.tex},
Theorem 1.1) states that $\Theta^\Phi_{R,\lambda}$ is $R$-positive as soon as
$\lambda\ge(R-1)\beta_2(\Phi)$, so the \emph{protected inclusion constant}
\begin{equation}
\label{eq:lamstar}
    \lambda_R^*(\Phi)
    :=\inf\{\lambda\ge0:\Theta^\Phi_{R,\lambda}\ \text{is $R$-positive}\}
\end{equation}
obeys $\lambda_R^*(\Phi)\le(R-1)\beta_2(\Phi)$.  Throughout, $R$-positivity is
used in its Choi form: $\Psi$ is $R$-positive if and only if
$\ip{\vecop C}{J(\Psi)\vecop C}\ge0$ for every $C$ with $\rank C\le R$, where
$J(\Psi)=\sum_{ij}E_{ij}\otimes\Psi(E_{ij})$.  In that normalization
$\Delta_d$ contributes $\norm C_F^2$ and $\id$ contributes $\abs{\Tr C}^2$, so
the correction in \cref{eq:theta} contributes
$\lambda(\norm C_F^2-\abs{\Tr C}^2/R)$.

The point of what follows is that the bound \cref{eq:lamstar} is attained by a
family whose coefficient system is far smaller than $M_d$.

\section{The edge-Kraus graph family}

Let $G=(V,E)$ be a simple graph on $V=\{1,\dots,m\}$ with $E\neq\emptyset$.
For an edge $e=\{i,j\}$ with $i<j$ set
\begin{equation}
\label{eq:kraus}
    L_e=E_{ij}-E_{ji}\in M_m,
    \qquad
    J=E_{01}-E_{10}\in M_2,
    \qquad
    A_e=L_e\otimes J\in M_{2m}.
\end{equation}
Both factors are real and skew, so $A_e^T=A_e=A_e^*$: the Kraus operators are
real symmetric, which is the frame-independent hypothesis pair amplification
needs.  Put
\begin{equation}
\label{eq:phig}
    \Phi_G(Y)=\sum_{e\in E}A_eYA_e .
\end{equation}

\begin{lemma}[Exact local constant]
\label{lem:beta}
$\beta_2(\Phi_G)=1$ for every nonempty $G$.
\end{lemma}

\begin{proof}
For $c=(c_e)_{e\in E}$ write $K_c=\sum_ec_eA_e=L_c\otimes J$ with
$L_c=\sum_ec_eL_e$.  The $L_e$ are Frobenius-orthogonal with
$\norm{L_e}_F^2=2$, so $\norm{L_c}_F^2=2\norm c_2^2$.  A complex skew matrix
has its nonzero singular values in equal pairs, hence
$2s_1(L_c)^2\le\norm{L_c}_F^2$, that is $s_1(L_c)^2\le\norm c_2^2$.  Since $J$
is unitary, the singular values of $L_c\otimes J$ are those of $L_c$ with
multiplicities doubled, so $s_1(K_c)=s_2(K_c)=s_1(L_c)$ and
\[
    s_1(K_c)^2+s_2(K_c)^2=2s_1(L_c)^2\le2\norm c_2^2 .
\]
A vector supported on a single edge attains equality, since $s_1(L_e)=1$.
\end{proof}

Write $\Theta^G_{R,\lambda}$ for \cref{eq:theta} with $\Phi=\Phi_G$, $d=2m$.

\begin{theorem}[Clique sharpness]
\label{thm:clique}
If $G$ contains an $R$-clique then $\lambda^*_R(\Phi_G)=R-1$.  More generally,
for every nonempty $S\subseteq V$ with $\abs S\le R$,
\begin{equation}
\label{eq:avgdeg}
    \lambda^*_R(\Phi_G)\ \ge\ \frac{2\abs{E(G[S])}}{\abs S}.
\end{equation}
\end{theorem}

\begin{proof}
The upper bound is pair amplification with \cref{lem:beta}.  For the lower
bound take $C=P_S\otimes E_{01}$, where $P_S=\sum_{i\in S}E_{ii}$.  Then
$\rank C=\abs S$, $\Tr C=0$ and $\norm C_F^2=\abs S$.  An edge $e$ inside $S$
contributes $-2$ to $\ip{\vecop C}{J(\Phi_G\circ\tau)\vecop C}$ and an edge
meeting $S$ in at most one vertex contributes $0$, so that pairing equals
$-2\abs{E(G[S])}$, while the correction contributes $\lambda\abs S$.
Nonnegativity forces \cref{eq:avgdeg}.  For an $R$-clique,
$2\abs{E(G[S])}=R(R-1)$ and the total is $R(\lambda-R+1)$.
\end{proof}

So the obstruction is the maximum induced average degree up to size $R$, and a
clique is exactly what makes the universal bound $R-1$ sharp.

\section{One layer is not enough}

The natural small coefficient system attached to $G$ is the graph operator
system
\begin{equation}
\label{eq:sg}
    \cS_G=\operatorname{span}\{E_{ii}:i\in V\}
      +\operatorname{span}\{E_{ij}:\{i,j\}\in E\}\subseteq M_m,
\end{equation}
and the map it carries is the graph reduction
$R_G(X)=\sum_eL_eX^TL_e^*$.  Let $p=\ketbra00$ and $q=\ketbra11$ in $M_2$.
Compressing the input to the $q$-corner and the output to the $p$-corner gives
\begin{equation}
\label{eq:corner}
    (I_m\otimes\bra0)\,(\Phi_G\circ\tau)(X\otimes q)\,(I_m\otimes\ket0)=R_G(X),
\end{equation}
which looks like the wanted compression.  It is not, because the same
compression annihilates the input:
$(I_m\otimes\bra0)(X\otimes q)(I_m\otimes\ket0)=0$.  Hence the identity term of
\cref{eq:theta} disappears and
\begin{equation}
\label{eq:corner-loss}
    \Delta_{2m}-\tfrac1R\id\ \longmapsto\ \Delta_m .
\end{equation}
The one-layer pencil therefore certifies the \emph{unprotected} hierarchy: for
regular $G$ containing an $s$-clique, $R_G+\lambda\Delta_m$ is $s$-positive
exactly for $\lambda\ge s-1$.  That is a true statement, and after unital
normalization an honest depolarizing mixture, but the protected direction --- the
whole content of pair amplification --- has been compressed away.

\section{The two-layer system}

Keep both diagonal qubit layers.  With $\cD_2=\operatorname{span}\{p,q\}$ put
\begin{equation}
\label{eq:ug}
    \cU_G:=\cS_G\otimes\cD_2
    =\{X\otimes p+Y\otimes q:X,Y\in\cS_G\}
    \cong\cS_G\oplus\cS_G,
\end{equation}
a unital self-adjoint operator system in $M_{2m}$.  Only \emph{diagonal} qubit
coefficients appear; the off-diagonal $E_{01}$ enters the witness below, never
the coefficient system.

\begin{proposition}[Invariance and layer action]
\label{prop:layers}
$\Phi_G\circ\tau$, $\Delta_{2m}$ and $\id$ preserve $\cU_G$, and
\begin{equation}
\label{eq:layer}
    (\Phi_G\circ\tau)(X\otimes p+Y\otimes q)
    =R_G(Y)\otimes p+R_G(X)\otimes q .
\end{equation}
Under $\cU_G\cong\cS_G\oplus\cS_G$,
\[
    \Theta^G_{R,\lambda}(X,Y)
    =\Bigl(R_G(Y)+\lambda(\Tr X+\Tr Y)I-\tfrac\lambda RX,\;
           R_G(X)+\lambda(\Tr X+\Tr Y)I-\tfrac\lambda RY\Bigr).
\]
\end{proposition}

\begin{proof}
Since $L_e^*=-L_e$ one has $\sum_eL_eX^TL_e=-R_G(X)$, and a direct computation
gives $JpJ=-q$ and $JqJ=-p$.  Hence
$A_e(X\otimes p)^TA_e=(L_eX^TL_e)\otimes(JpJ)$, and the two signs cancel,
leaving $R_G(X)\otimes q$; likewise for the $q$-layer.  The correction terms
are immediate.
\end{proof}

The construction is equivariant under $\operatorname{Aut}(G)\times\mathbb Z_2$,
the second factor exchanging layers.  If $G$ is $k$-regular then
$R_G(I)=kI$, because $L_e^2=-(E_{ii}+E_{jj})$ and each vertex lies in $k$
edges; consequently
\begin{equation}
\label{eq:unital}
    \Theta^G_{R,\lambda}(I_{2m})=D_{G,R,\lambda}I_{2m},
    \qquad
    D_{G,R,\lambda}=k+\lambda\Bigl(2m-\tfrac1R\Bigr),
\end{equation}
so $T^G_{R,\lambda}:=D_{G,R,\lambda}^{-1}\Theta^G_{R,\lambda}|_{\cU_G}$ is
unital and Hermiticity-preserving.  For $k$-regular $G$,
\begin{equation}
\label{eq:dim}
    \dim_\C\cS_G=m+2\abs E=m(1+k),
    \qquad
    \dim_\C\cU_G=2m(1+k),
\end{equation}
so a monic pencil on a real basis of the traceless self-adjoint part of $\cU_G$
uses $g=2m(1+k)-1$ variables rather than $(2m)^2-1$.

\section{The inclusion dictionary}

The bridge from positivity levels to spectrahedral inclusion is standard in the
matrix-convexity literature (Davidson--Dor-On--Shalit--Solel,
\href{https://arxiv.org/abs/1601.07993}{arXiv:1601.07993}); we record the form
used here, which does not require a minimal pencil.

\begin{lemma}[Inclusion dictionary]
\label{lem:dict}
Let $\cU\subseteq M_d$ be an operator system with $F_1,\dots,F_g$ a real basis
of its traceless self-adjoint part, and let $T:\cU\to M_d$ be unital and
Hermiticity-preserving.  With
$L_F(X)=I_d\otimes I_n+\sum_aF_a\otimes X_a$ and
$L_B(X)=I_d\otimes I_n+\sum_aT(F_a)\otimes X_a$, and $\cD_F(n),\cD_B(n)$ their
positivity loci,
\[
    \cD_F(n)\subseteq\cD_B(n)
    \quad\Longleftrightarrow\quad
    T\ \text{is $n$-positive on }\cU .
\]
\end{lemma}

\begin{proof}
If $T$ is $n$-positive then $L_B(X)=(T\otimes\id_n)(L_F(X))$ and inclusion is
immediate.  Conversely let $Y\in M_n(\cU)_+$, with its unique expansion
$Y=I_d\otimes Y_0+\sum_aF_a\otimes Y_a$, where $Y_0=d^{-1}\Tr_{M_d}Y\succeq0$.
If $Y_0$ is invertible, conjugating by $I_d\otimes Y_0^{-1/2}$ places $Y$ in
the monic slice parametrized by $L_F$; inclusion and conjugating back give
$(T\otimes\id_n)(Y)\succeq0$.  For singular $Y_0$ apply this to
$Y+\varepsilon I_d\otimes I_n$ and let $\varepsilon\downarrow0$.
\end{proof}

\section{The witness survives the compression}

This is the step the one-layer system fails.  Let $S=\{s_1,\dots,s_s\}$ be an
$s$-clique, let $e_1,\dots,e_s$ be the standard basis of the level space
$\C^s$, and let $f_i$ be the vertex basis of $\C^m$.  Put
\begin{equation}
\label{eq:witness}
    \xi_S=\sum_{a=1}^se_a\otimes(f_{s_a}\otimes\ket1),
    \qquad
    \eta_S=\sum_{a=1}^se_a\otimes(f_{s_a}\otimes\ket0),
\end{equation}
so that
$Y_S=\ketbra{\xi_S}{\xi_S}=\sum_{a,b}E_{ab}\otimes(E_{s_as_b}\otimes q)$
lies in $M_s(\cU_G)_+$: it is positive of rank one, and it lies in
$M_s(\cU_G)$ because a clique gives $P_S\cS_GP_S=M_s$.

\begin{lemma}[Compressed clique expectation]
\label{lem:witness}
For every $R\ge1$ and $\lambda\ge0$,
\begin{equation}
\label{eq:witness-val}
    \ip{\eta_S}{(\id_s\otimes\Theta^G_{R,\lambda})(Y_S)\,\eta_S}
    =s\bigl(\lambda-(s-1)\bigr).
\end{equation}
\end{lemma}

\begin{proof}
By \cref{eq:layer} each ordered pair $a\neq b$ inside the clique contributes
$-1$ and diagonal pairs contribute $0$, so the $\Phi_G\circ\tau$ term gives
$-s(s-1)$.  The $\Delta_{2m}$ term gives $\lambda s$.  The identity term sends
the input $q$-layer to the $q$-layer, whereas $\eta_S$ lies in the orthogonal
$p$-layer, so the protected term $-\lambda R^{-1}\id$ contributes nothing.
\end{proof}

Note what \cref{eq:witness-val} does \emph{not} depend on: the value is
independent of $R$.  The off-diagonal $E_{01}$ has been retained exactly where
it is needed, as the transition between $\xi_S$ and $\eta_S$; it never had to
enter the coefficient system.  Equivalently, \cref{lem:witness} is the rank-$s$
Choi witness $C=P_S\otimes E_{01}$ of \cref{thm:clique}, read after
compression.

\begin{theorem}[Exact protected inclusion on $\cU_G$]
\label{thm:inclusion}
Let $G$ be $k$-regular and contain an $R$-clique, and let
$B^{(R,\lambda)}_a=T^G_{R,\lambda}(F_a)$ for a traceless self-adjoint basis
$\{F_a\}$ of $\cU_G$.  Then
\begin{equation}
\label{eq:exact}
    \inf\bigl\{\lambda\ge0:\cD_F(R)\subseteq\cD_{B^{(R,\lambda)}}(R)\bigr\}
    =R-1 .
\end{equation}
If moreover $G$ contains an $(R+1)$-clique, then at $\lambda=R-1$
\begin{equation}
\label{eq:levels}
    \cD_F(n)\subseteq\cD_{B^{(R,R-1)}}(n)
    \quad\Longleftrightarrow\quad
    n\le R .
\end{equation}
\end{theorem}

\begin{proof}
By \cref{lem:beta} and pair amplification, $\Theta^G_{R,R-1}$ is $R$-positive on
all of $M_{2m}$, hence on $\cU_G$; raising $\lambda$ adds an $R$-positive map.
For $\lambda<R-1$, \cref{lem:witness} with $s=R$ exhibits a positive element of
$M_R(\cU_G)$ with negative image.  \Cref{lem:dict} converts both statements
into \cref{eq:exact}.  At $\lambda=R-1$, $R$-positivity gives $n$-positivity for
$n\le R$; an $(R+1)$-clique and \cref{lem:witness} with $s=R+1$ give the value
$-(R+1)<0$, so $(R+1)$-positivity fails, and positivity levels are nested.
\end{proof}

\section{The rook graph}

Let $G=H(2,q)$ on $V=[q]\times[q]$, two vertices adjacent when they agree in
exactly one coordinate.  Then $m=q^2$, $k=2(q-1)$, the clique number is $q$
(rows and columns), and the adjacency eigenvalues are $2(q-1)$, $q-2$, $-2$.

\begin{corollary}
\label{cor:rook}
For every $1\le R\le q$, \cref{eq:exact} holds with $\lambda^*=R-1$, and for
$1\le R<q$ the exact level separation \cref{eq:levels} holds.  The protected
pencil uses
\[
    g_{\mathrm{prot}}(q)=2q^2(2q-1)-1
\]
traceless coefficients, against $4q^4-1$ for a tomographically complete
dictionary on $M_{2q^2}$.
\end{corollary}

On $\cS_G$ the graph reduction $R_G$ acts with eigenvalue $k=2(q-1)$ on the
identity and $\mu\in\{q-2,-2,-1\}$ on the traceless part, the value $-1$ being
the edge-coefficient eigenvalue.  Writing
$D^{\mathrm{prot}}_{q,R,\lambda}=2(q-1)+\lambda(2q^2-R^{-1})$, a traceless
$X\in\cS_G$ with $R_G(X)=\mu X$ gives layer-even and layer-odd eigenvectors
\[
    T^G_{R,\lambda}(X,X)=\frac{\mu-\lambda/R}{D^{\mathrm{prot}}_{q,R,\lambda}}(X,X),
    \qquad
    T^G_{R,\lambda}(X,-X)=-\frac{\mu+\lambda/R}{D^{\mathrm{prot}}_{q,R,\lambda}}(X,-X),
\]
so the whole protected pencil is governed by a handful of association-scheme
and layer-parity scalars.  The normalized map is the convex combination
\[
    T^G_{R,\lambda}
    =(1-p)\,\frac{\Phi_G\circ\tau}{k}
     +p\,\frac{\Delta_{2m}-R^{-1}\id}{2m-R^{-1}},
    \qquad
    p=\frac{\lambda(2m-R^{-1})}{k+\lambda(2m-R^{-1})} .
\]

\begin{remark}[The weight is not a noise probability]
At $\lambda=R-1$ the weight is
$p^{\mathrm{prot}}_R(q)=(R-1)(2q^2-R^{-1})\big/\bigl(2(q-1)+(R-1)(2q^2-R^{-1})\bigr)$.
This is a protected-ray weight, not a depolarizing probability: the second
summand is $R$-positive but not completely positive.  The one-layer corner of
\cref{eq:corner-loss} \emph{does} give a genuine depolarizing mixture, with
threshold $p^*_R(q)=q^2(R-1)/\bigl(2(q-1)+q^2(R-1)\bigr)$ and
$g_{\mathrm{corner}}(q)=q^2(2q-1)-1$ coefficients.  The two should not be
conflated: the smaller count certifies the unprotected corner.
\end{remark}

\section*{What is and is not established}

Established: the exact constant $\beta_2(\Phi_G)=1$; attainment of the
amplified coefficient $R-1$ for graphs with an $R$-clique, and the induced
average-degree lower bound \cref{eq:avgdeg}; invariance of
$\cU_G=\cS_G\otimes\cD_2$ with the layer action \cref{eq:layer}; a positive
element of $M_R(\cU_G)$ witnessing sharpness after compression; and exact
inclusion through level $R$ with failure at $R+1$.  The formal-verification
boundary is recorded below.

Not established.  Minimality of $\cU_G$ among operator systems retaining the
witness is not claimed.  Neither is an exact graph invariant replacing clique
number in general: \cref{eq:avgdeg} is a lower bound and $R-1$ an upper bound,
and closing that gap for clique-free graphs is a separate extremal problem.
The Bose--Mesner algebra alone does not suffice, being commutative, so that
positivity of a map out of it already forces complete positivity.  Finally, we
have not carried out a literature priority check for this formulation; targeted
searches did not surface it, which is not the same thing.

\section*{Formal verification}

\Cref{lem:beta}, both halves of \cref{thm:clique}, and the finite-coordinate
content of \cref{prop:layers} and \cref{lem:witness} are machine-checked in the
Lean~4 development accompanying \texttt{rank-r.tex},
\url{https://github.com/elazarg/rank-r}, in the files \texttt{GraphFamily.lean},
\texttt{Theta.lean}, \texttt{ThetaBound.lean} and \texttt{GraphLayers.lean}.

\begin{center}
\begin{tabular}{ll}
\hline
Here & Lean\\
\hline
$A_e$ of \cref{eq:kraus} & \texttt{edgeKraus}, \texttt{graphKraus}\\
\cref{lem:beta} & \texttt{choiTwoBound\_graphKraus\_iff}\\
$R$-positivity of \cref{eq:theta} & \texttt{ThetaPositive}\\
$\lambda^*_R(\Phi)\le(R-1)\beta_2(\Phi)$ of \cref{eq:lamstar}
  & \texttt{thetaPositive\_of\_choiTwoBound}\\
\cref{eq:avgdeg} & \texttt{two\_mul\_card\_le\_of\_thetaPositive}\\
\cref{thm:clique}, clique case & \texttt{thetaPositive\_graphKraus\_iff}\\
\cref{prop:layers} & \texttt{graphTheta\_layers},
  \texttt{graphTheta\_mem\_graphTwoLayerSpace}\\
\cref{lem:witness} & \texttt{cliqueLayerWitness\_posSemidef},
  \texttt{cliqueLayerWitness\_mem\_matrixLevelSpace},
  \texttt{qform\_graphThetaCliqueOutput}\\
\hline
\end{tabular}
\end{center}

Two of those are stronger than stated here.  \Cref{lem:beta} is formalized
two-sidedly --- the constants valid for $\Phi_G$ are exactly the reals $\ge1$ ---
so $\beta_2(\Phi_G)=1$ is certified as an equality rather than as a bound; and
the clique case of \cref{thm:clique} is an equivalence, $\Theta^G_{R,\lambda}$
being $R$-positive precisely for $\lambda\ge R-1$.

The formal upper bound does not follow the singular-value argument given above.
Every $A_e$ lies in the double-skew subspace, so the antisymmetrizer $Q_-$ fixes
its vectorization and, being self-adjoint, lets each pairing
$\ip{\operatorname{vec}A_e}{\operatorname{vec}M}$ be read at
$Q_-\operatorname{vec}M$ instead; the $\operatorname{vec}A_e$ are orthogonal of
squared norm $4$, so Bessel reduces the estimate to
$\norm{Q_-}_{S(2)}=\frac12$, which \texttt{rank-r.tex} already carries.  The
lower bound runs on the witness $C=P_S\otimes E_{01}$, which in the development
is the extremizer already built for the sharpness half of Theorem~1.1 there.

Not formalized: \cref{lem:dict} and \cref{thm:inclusion}, together with the
general free-spectrahedral dictionary they rest on.  The concrete spaces
$\cS_G$ and $\cU_G$, their unital and adjoint structure, the graph-reduction
range, the displayed layer formula, invariance of the corrected map, and the
positive compressed clique input with its exact expectation are formalized.
The remaining statements restate \cref{thm:clique} in matrix-convexity language
and introduce no new constant.

\section*{Numerical corroboration}

Every claim above was checked numerically by building $\Phi_G$, $\Theta^G$,
$\cU_G$ and the witnesses explicitly.  For $G$ ranging over $K_4$ with a
pendant, $C_5$, the Petersen graph and $H(2,3)$, in every case to
machine zero: the bound $s_1(K_c)^2+s_2(K_c)^2\le2$ over $2\cdot10^4$ random
unit $c$, together with equality at each single-edge coefficient vector, so that
$\beta_2(\Phi_G)=1$ is attained and not merely approached; the Choi form at
$C=P_S\otimes E_{01}$ against every clique; the layer action \cref{eq:layer};
$R_G(I)=kI$ for regular $G$; the compressed expectation \cref{eq:witness-val},
including for $R\neq s$, confirming that its value does not depend on $R$; and
that $Y_S$ is positive of rank one.  For $H(2,q)$ with $q\in\{3,4\}$ it confirms
$\dim\cS_G=m(1+k)$ and that the spectrum of $R_G$ on $\cS_G$ is
$\{2(q-1)\}\cup\{q-2,-2,-1\}$.

The extremizer of \cref{lem:beta} is a coordinate direction, so a random search
over $c$ converges to a strict underestimate --- $0.910$ for $H(2,3)$ --- and
would test the wrong thing.  What is checked is the upper bound over random $c$
\emph{together with} attainment at $c=\delta_e$.

\end{document}
